% Relatorio: fronteira imersa com forcamento direto no HiGFlow
%
% Compila com pdflatex + bibtex + pdflatex + pdflatex, ou latexmk -pdf main
% No Overleaf: enviar a pasta inteira; o arquivo principal e' este.
%
% Os NUMEROS vivem em dados/*.dat e sao lidos por pgfplotstable, nao digitados
% no texto.  Regenerar uma medida troca o .dat e o relatorio acompanha -- sem
% a chance de a tabela e o texto discordarem, que e' o erro classico aqui.
\documentclass[11pt,a4paper]{article}

\usepackage[T1]{fontenc}
\usepackage[utf8]{inputenc}
\usepackage[brazil]{babel}
\usepackage{lmodern}
\usepackage{microtype}
\usepackage{amsmath,amssymb}
\usepackage{booktabs}
\usepackage{siunitx}
\usepackage{graphicx}
\usepackage{pgfplots}
\usepackage{pgfplotstable}
\usepackage{listings}
\usepackage{xcolor}
\usepackage[hidelinks]{hyperref}
\usepackage[margin=2.7cm]{geometry}

\pgfplotsset{compat=1.18}
\sisetup{output-decimal-marker={,}, group-separator={.}}

\lstset{
  basicstyle=\ttfamily\footnotesize,
  keywordstyle=\color{teal},
  commentstyle=\color{gray},
  breaklines=true,
  frame=leftline,
  framesep=6pt,
  xleftmargin=10pt,
}

\newcommand{\Cd}{C_{\mathrm{d}}}
\newcommand{\Cl}{C_{\mathrm{l}}}
\newcommand{\dt}{\Delta t}

\title{Fronteira imersa com forçamento direto no HiGFlow:\\
       implementação, verificação e o experimento de Schäfer--Turek}
\author{Antonio Castelo Filho}
\date{\today}

\begin{document}
\maketitle

\begin{abstract}
Este relatório descreve a implementação de um método de fronteira imersa com
forçamento direto no HiGFlow, sua verificação, e o estudo de refinamento sobre o
escoamento em canal com obstáculo cilíndrico (\emph{benchmark} DFG/Schäfer--Turek
2D-1, $\mathrm{Re}=20$). Duas formulações foram implementadas e comparadas: a
força \emph{defasada}, calculada com a velocidade do passo anterior, e a rota
\emph{pós-preditor} de Fadlun/Uhlmann, com forçamento iterado. Registra-se
também, deliberadamente, o que cada medida \emph{não} estabelece --- incluindo
três oráculos que passaram com o método errado.
\end{abstract}

\tableofcontents
\bigskip

\section{Objetivo e escopo}

O HiGFlow resolve Navier--Stokes incompressível sobre malha cartesiana em
octree (HiGTree). Corpos sólidos de geometria curva não são representáveis por
essa malha sem escada de células, e é esse o problema que a fronteira imersa
resolve: o corpo é descrito por uma malha \emph{lagrangeana} de marcadores que
não coincide com a malha \emph{euleriana} do escoamento, e o acoplamento entre
as duas se faz por um campo de força.

Este relatório cobre:

\begin{itemize}
  \item a formulação implementada e as duas variantes comparadas
        (Seção~\ref{sec:formulacao});
  \item as decisões de implementação e estrutura de dados
        (Seção~\ref{sec:implementacao});
  \item a verificação dos operadores de transferência, isolada do solver
        (Seção~\ref{sec:verificacao});
  \item o experimento --- canal com obstáculo cilíndrico, \emph{benchmark}
        DFG/Schäfer--Turek 2D-1 --- e o estudo de refinamento
        (Seções~\ref{sec:experimento} e~\ref{sec:resultados}).
\end{itemize}

\paragraph{Escopo deliberadamente restrito.} As decisões abaixo foram tomadas no
início e limitam o que este trabalho afirma:

\begin{center}
\begin{tabular}{ll}
  \toprule
  corpo    & rígido e \textbf{fixo} ($U_{\mathrm{corpo}} = 0$) \\
  alcance  & apenas o solver newtoniano \\
  malha    & uniforme (sem refinamento local) \\
  \bottomrule
\end{tabular}
\end{center}

Corpo móvel, corpo deformável e interface entre fluidos \emph{não} são
extensões diretas deste trabalho; a Seção~\ref{sec:discussao} explica por quê.

\paragraph{Uma nota sobre o que este relatório registra.} Além dos resultados,
registram-se as medidas que \emph{não} discriminaram --- oráculos que passaram
com o método errado --- porque elas são a parte transferível da experiência. Um
resultado correto obtido por um teste que não podia detectar o erro é sorte, e
sorte não se reproduz.

\section{Formulação}
\label{sec:formulacao}

\subsection{Os dois operadores de transferência}

A fronteira é discretizada em marcadores $\mathbf{X}_k$, $k = 1,\dots,N_L$, cada
um com uma medida geométrica $w_k$ (comprimento de arco em 2D, área em 3D). O
acoplamento entre as malhas se faz por dois operadores, construídos sobre o mesmo
núcleo regularizado $\delta_h$:

\begin{align}
  \mathbf{u}(\mathbf{X}_k) &= \sum_{\mathbf{x}} \mathbf{u}(\mathbf{x})\,
      \delta_h(\mathbf{x}-\mathbf{X}_k)\, h^{D}
      && \text{(interpolação)} \label{eq:interp}\\
  \mathbf{F}(\mathbf{x}) &= \sum_{k} \mathbf{f}_k\,
      \delta_h(\mathbf{x}-\mathbf{X}_k)\, \Delta V_k
      && \text{(espalhamento)} \label{eq:espalha}
\end{align}

onde $D$ é a dimensão do espaço e $\Delta V_k$ é o \emph{volume} do marcador.

O núcleo é o de três pontos de Roma, Peskin e Berger~\cite{roma1999}:
\begin{equation}
  \phi(r) =
  \begin{cases}
    \tfrac{1}{3}\left(1 + \sqrt{1-3r^2}\right), & |r| \le \tfrac{1}{2},\\[4pt]
    \tfrac{1}{6}\left(5 - 3|r| - \sqrt{1-3(1-|r|)^2}\right),
      & \tfrac{1}{2} \le |r| \le \tfrac{3}{2},\\[4pt]
    0, & \text{caso contrário,}
  \end{cases}
  \qquad
  \delta_h(\mathbf{x}) = \frac{1}{h^{D}}\prod_{d=1}^{D}\phi\!\left(\frac{x_d}{h}\right).
\end{equation}

\paragraph{O par adjunto é o que importa.} Com o \emph{mesmo} núcleo nos dois
operadores, espalhar é o transposto de interpolar, e disso decorre a conservação:
a força total que sai dos marcadores chega às facetas,
\begin{equation}
  \sum_{\mathbf{x}} \mathbf{F}(\mathbf{x})\,h^{D}
  \;=\; \sum_k \mathbf{f}_k \,\Delta V_k .
  \label{eq:conservacao}
\end{equation}
Trocar um dos dois núcleos quebra~\eqref{eq:conservacao} sem produzir nenhum
sintoma visível --- ponto ao qual a Seção~\ref{sec:verificacao} retorna.

\paragraph{O volume do marcador não é a medida geométrica.} Para um corpo de
codimensão~1 (curva em 2D, superfície em 3D), o marcador representa um retalho de
espessura uma célula:
\begin{equation}
  \Delta V_k = w_k \, h^{\,\mathrm{codim}} , \qquad \mathrm{codim} = 1 .
\end{equation}
O expoente é a \textbf{codimensão do corpo}, não a dimensão do espaço. Em 2D as
duas coincidem, e usar $h^{D-1}$ passa despercebido; em 3D isso daria
$\mathrm{d}A\cdot h^2$ em vez de $\mathrm{d}A\cdot h$, errando a força por um
fator $1/h$.

\subsection{Forçamento direto: duas variantes}

Para corpo rígido e fixo, a força não tem lei constitutiva --- é um
multiplicador de Lagrange que vale o que for preciso para impor
$\mathbf{u}=\mathbf{0}$ na fronteira:
\begin{equation}
  \mathbf{f}_k = \frac{\mathbf{0} - \mathbf{u}(\mathbf{X}_k)}{\dt}.
  \label{eq:forca}
\end{equation}

A diferença entre as duas variantes implementadas está em \emph{de onde vem}
$\mathbf{u}(\mathbf{X}_k)$ em~\eqref{eq:forca}, e \emph{quando} a força é
aplicada.

\paragraph{Defasada.} A velocidade é a do passo anterior, $\mathbf{u}^n$, e a
força é somada ao campo de fonte \emph{antes} do preditor. Custa uma chamada
num único ponto do passo, e impõe o não escorregamento a $O(\dt)$.

\paragraph{Pós-preditor (Fadlun/Uhlmann~\cite{uhlmann2005}).} O preditor produz
a velocidade provisória $\mathbf{u}^*$ \emph{sem} força; a força é então
calculada de $\mathbf{u}^*$ e corrige a própria $\mathbf{u}^*$:
\begin{equation}
  \mathbf{u}^* \leftarrow \mathbf{u}^* + \alpha\,\dt\,\mathbf{F}, \qquad
  \mathbf{f}_k = \frac{\mathbf{0}-\mathbf{u}^*(\mathbf{X}_k)}{\dt}.
\end{equation}
O termo $O(\dt)$ da defasagem desaparece. A projeção de pressão roda
\emph{depois} e desfaz parte da correção, de modo que o não escorregamento
continua aproximado.

\paragraph{Forçamento iterado.} Uma aplicação não zera o escorregamento porque
espalhar e interpolar não comutam ($\mathcal{S}\mathcal{I} \neq \mathcal{I}$).
Repetir interpolação--força--espalhamento dentro do passo é o
\emph{multi-direct forcing}~\cite{wang2008}; a literatura relata de~5 a~20
iterações.

\subsection{O que ``segunda ordem'' significa na literatura}

O método de fronteira imersa formalmente de segunda ordem é o de Griffith
\emph{et al.}~\cite{griffith2007}, cuja base não adaptativa é Lai e
Peskin~\cite{lai2000}. A distinção importa e é frequentemente perdida: a segunda
ordem é \textbf{do esquema}, não da imposição da condição de contorno. Nas
palavras de~\cite{lai2000}, o delta discreto ``impede o método de fronteira
imersa de ser mais que de primeira ordem'', porque a velocidade não tem derivada
contínua através da fronteira; o ganho efetivo é \emph{menos viscosidade
numérica}. Ceniceros \emph{et al.}~\cite{ceniceros2010} medem exatamente isso:
ordem~1 na vizinhança da interface, tendendo a~2 longe dela.

\section{Implementação}
\label{sec:implementacao}

\subsection{Onde o código vive}

\begin{center}
\begin{tabular}{ll}
  \toprule
  \texttt{higflow/src/hig-flow-fronteira-imersa.\{h,c\}} & o núcleo do método \\
  \texttt{higflow/examples-common/fronteira-imersa.c}    & adaptador solver/módulo \\
  \texttt{higflow/fronteira-imersa/testes/}              & verificação isolada \\
  \texttt{higflow/example2d\_SchaeferTurek/}             & o experimento \\
  \bottomrule
\end{tabular}
\end{center}

\paragraph{O módulo não conhece o solver, e isso é deliberado.} Ele recebe um
\texttt{sim\_facet\_domain} e uma \texttt{distributed\_property} por direção, não
um \texttt{higflow\_solver}. Foi essa separação que permitiu verificar a
conservação sem montar um solver (Seção~\ref{sec:verificacao}).

\paragraph{O acoplamento é por ponteiro de função.} O arquivo
\texttt{hig-flow-step.c} é ligado por \emph{todos} os exemplos. Uma chamada
direta obrigaria os doze \emph{Makefiles} a ligar o módulo por um recurso que um
exemplo usa. Com o gancho, quem não instala corpo não referencia símbolo algum.
A primeira tentativa, com chamada direta, levou a suíte de 33 para 0 de 33 casos.

\subsection{Estrutura de dados da malha lagrangeana}

Os marcadores vivem num \texttt{DMSwarm} do PETSc --- biblioteca já presente no
projeto, sem dependência nova. A entrada é uma \textbf{curva fechada por
segmentos}; em 3D, a extrusão dessa curva, formando uma superfície.

\paragraph{Marcadores nos centros dos subsegmentos, não nos vértices.} Marcador
em vértice duplicaria o vértice compartilhado entre segmentos consecutivos, e a
soma dos pesos deixaria de ser o perímetro. Com centros, o peso de cada marcador
é o comprimento do seu subsegmento, e a soma dá o perímetro exato --- invariante
que o teste cobra com tolerância de $10^{-12}$.

\paragraph{Distribuição por posse euleriana.} Cada marcador pertence ao rank que
possui a célula que o contém, condição para o espalhamento ser local. Uma guarda
aborta se a soma dos marcadores reivindicados não bater com o total: marcador
sobre fronteira de partição reivindicado por dois ranks (ou por nenhum) daria
força errada sem sintoma visível.

\subsection{O suporte do marcador, sem varredura}

O núcleo tem suporte de $1{,}5$ células para cada lado. Localizar as facetas do
suporte varrendo seria $O(N_L \times N_{\mathrm{facetas}})$. Não é preciso: a
malha é uniforme e os centros do suporte estão em \emph{deslocamentos conhecidos}
a partir de uma âncora, cada um achado por uma localização de ponto.

\paragraph{A âncora sai da célula, não de uma faceta.} A função
\texttt{sfd\_get\_facet\_with\_point} acha a faceta cujo \emph{plano} contém o
ponto; um marcador no meio de uma célula não está sobre plano algum naquela
direção. Ancorar numa faceta funcionava apenas para a direção em que o marcador
por acaso caísse sobre um plano de face --- e numa curva alinhada com a malha
isso atinge metade dos marcadores, com a força saindo exatamente pela metade.

\subsection{Acumulação franja \texorpdfstring{$\to$}{->} dono}

O espalhamento de um marcador próximo da borda escreve em facetas de \emph{outro}
rank. A sincronização da HiGTree (\texttt{dp\_sync}) envia dono\,$\to$\,franja e
\textbf{sobrescreve}: contribuições escritas numa franja seriam descartadas em
silêncio. A direção que falta é feita com \texttt{PetscSFReduce} e
\texttt{MPI\_SUM}, sobre um grafo montado a partir do que a HiGTree já tem
(\texttt{gid\_map}, \texttt{local\_count}, \texttt{total\_count}).

\paragraph{Contrato: o espalhamento é dono do campo que recebe.} Ele acumula
franja\,$\to$\,dono, o que só é correto se o campo contiver \emph{apenas}
contribuições espalhadas. Passar um campo já povoado faz o \emph{reduce} somar
um campo inteiro nas fronteiras de partição. Para somar num campo povoado ---
como $\mathbf{u}^*$ na rota pós-preditor --- espalha-se num campo zerado próprio
e soma-se depois.

\section{Verificação}
\label{sec:verificacao}

Os operadores de transferência são verificados \emph{sem solver}, em domínio
pequeno, com cinco afirmações em $np = 1, 2, 3$:

\begin{enumerate}
  \item \textbf{Peso total igual à medida geométrica.} A curva é um quadrado de
        lado $0{,}5$: perímetro exatamente $2{,}0$. Círculo não serviria --- o
        polígono inscrito tem perímetro \emph{próximo}, e ``próximo'' não
        distingue subdivisão errada de discretização.
  \item \textbf{Interpolar campo constante devolve a constante.} É a partição da
        unidade do núcleo. Pega núcleo errado, suporte incompleto e normalização
        errada.
  \item \textbf{Conservação da força}, equação~\eqref{eq:conservacao}. É o único
        dos três que exercita a acumulação franja\,$\to$\,dono.
\end{enumerate}

\subsection{Três oráculos que não discriminaram}

Esta subseção é o registro do que \emph{não} funcionou como teste, e é a parte
transferível da experiência.

\paragraph{A conservação da força não pega erro nos marcadores.} A
equação~\eqref{eq:conservacao} é uma \textbf{identidade em $\Delta V_k$}: ela é
verdadeira seja qual for o significado do peso, e os dois lados usam os mesmos
marcadores. Ela passou, exata, em duas situações de erro grave: com metade dos
marcadores devolvendo zero (âncora na faceta), e com a força vinte vezes maior
que o devido (peso $=\mathrm{d}s$ em vez de $\mathrm{d}s\cdot h$). Quem pegou os
dois foi a afirmação~2, que compara com um valor absoluto \emph{conhecido}.

\paragraph{A esteira do cilindro não é oráculo.} Ela parece correta numa faixa
larga de implementações erradas. O canal com obstáculo \emph{é} bom caso ---
desde que se compare número com número contra valores publicados, não figura.

\paragraph{Taxa de convergência com número de passos fixo não mede taxa.}
Refinar $\dt$ mantendo \texttt{numsteps} faz cada corrida terminar num
\textbf{tempo físico diferente}: comparam-se pontos distintos do transiente. Para
acurácia, o tempo final deve ser fixo. (Para \emph{estabilidade} vale o oposto:
instabilidade de realimentação cresce por passo, e ali o número de passos é que
deve ser fixo.)

\paragraph{Corolário de método.} Em três ocasiões distintas um contador ou uma
métrica repetiu \emph{exatamente} o mesmo valor após uma correção, e isso foi
lido como ``correção que errou por pouco''. Em todas, o significado era outro:
\textbf{ramo não exercitado}. Número que não muda quando deveria mudar é
informação, não ruído.

\section{O experimento}
\label{sec:experimento}

\subsection{Geometria e escala}

O caso é o \emph{benchmark} DFG/Schäfer--Turek 2D-1~\cite{schafer1996}:
escoamento laminar estacionário em canal sobre cilindro, $\mathrm{Re}=20$. Na
forma original:

\begin{center}
\begin{tabular}{ll}
  \toprule
  canal    & $[0;\,2{,}2] \times [0;\,0{,}41]$ \\
  cilindro & $D = 0{,}1$, centro em $(0{,}2;\,0{,}2)$ \\
  entrada  & $u(0,y) = 4\,U_m\,y(H-y)/H^2$, \; $U_m = 0{,}3$ \\
  fluido   & $\nu = 10^{-3}$, $\rho = 1$ \\
  \bottomrule
\end{tabular}
\end{center}

com $\bar{u} = 2U_m/3 = 0{,}2$ e $\mathrm{Re} = \bar{u}D/\nu = 20$.

\paragraph{Adimensionalização, e por que ela não é opcional.} O HiGFlow aplica
$1/\mathrm{Re}$ ao termo viscoso. Fornecer a geometria dimensional com
$\mathrm{Re}=20$ produziria um Reynolds efetivo de $0{,}4$. Todas as medidas
abaixo usam a escala de $D$ e da velocidade média:

\begin{center}
\begin{tabular}{ll}
  \toprule
  canal    & $22 \times 4{,}1$ \\
  cilindro & $D = 1$, centro em $(2;\,2)$ \\
  entrada  & $u(y) = 6\,y(4{,}1-y)/4{,}1^2$, \; média $1$, máximo $1{,}5$ \\
  $\mathrm{Re}$ & $20$ \\
  \bottomrule
\end{tabular}
\end{center}

$\Cd$ e $\Cl$ são adimensionais e não mudam com essa escala.

\paragraph{O cilindro está fora do eixo de propósito.} O centro em $y=2$ num
canal $[0;4{,}1]$ está a $2$ do fundo e $2{,}1$ do topo. A assimetria é do
\emph{benchmark} e é ela que produz o $\Cl$ publicado; centralizar invalidaria a
comparação.

\subsection{Coeficientes e valores de referência}

Nesta escala ($\rho=1$, $\bar{u}=1$, $D=1$),
\begin{equation}
  \Cd = -2\,F_x, \qquad \Cl = -2\,F_y,
\end{equation}
com $\mathbf{F} = \sum_k \mathbf{f}_k \Delta V_k$ a força que o corpo aplica ao
\emph{fluido}; a força hidrodinâmica sobre o corpo é a reação, daí o sinal.

\paragraph{Com forçamento iterado, a força do passo é a soma sobre as
iterações.} Cada iteração acrescenta uma parcela, e a última é a menor; ler
apenas a última daria arrasto uma ordem de grandeza pequeno demais.

Os valores de referência, de Nabh~\cite{nabh1998} e confirmados
por~\cite{john2001}:
\begin{center}
\begin{tabular}{lS[table-format=1.11]}
  \toprule
  $\Cd$       & 5.57953523384 \\
  $\Cl$       & 0.010618948146 \\
  $\Delta p$  & 0.11752016697 \\
  \bottomrule
\end{tabular}
\end{center}

\paragraph{$\Cl$ não é critério de acurácia aqui.} O deslocamento do cilindro
fora do eixo é $D/20$, e o suporte do núcleo de três pontos é $3h$. Com
$h = D/20$, o suporte cobre $0{,}15\,D$ --- \textbf{três vezes maior que a
própria assimetria} que produz o $\Cl$. Ele serve de diagnóstico de simetria,
não de critério.

\section{Resultados}
\label{sec:resultados}

Salvo indicação, as medidas usam $h = 0{,}05$ (20 células por diâmetro), $np=3$,
e o resíduo de não escorregamento é $\max_k |\mathbf{u}(\mathbf{X}_k)|$, a ser
comparado com $\max|\mathbf{u}| \approx 1{,}8$ no domínio.

\subsection{Estabilidade: o maior \texorpdfstring{$\dt$}{dt}, forçamento defasado}

\pgfplotstabletypeset[
  % TODAS string: os valores divergidos estouram o parser numerico do TeX.
  columns/0/.style={column name={$\dt$}, string type},
  columns/1/.style={column name={resíduo}, string type},
  columns/2/.style={column name={$\max|\mathbf{u}|$}, string type},
  columns/3/.style={column name={veredito}, string type},
  every head row/.style={before row=\toprule, after row=\midrule},
  every last row/.style={after row=\bottomrule},
]{dados/varre-dt-defasado.dat}

\medskip
Número de passos \emph{fixo} (400) para todos os valores: instabilidade de
realimentação cresce por passo, e a tempo fixo cada $\dt$ teria um número
diferente de oportunidades de amplificação.

\paragraph{A conclusão é o contrário do que motivou o experimento.} O
despenhadeiro está entre $5\times 10^{-2}$ e $10^{-1}$, \emph{acima} do CFL
convectivo ($\approx 0{,}026$): a fronteira imersa não é o gargalo de
estabilidade. O limite útil é de \textbf{acurácia} --- em $\dt = 5\times10^{-2}$,
último estável, o resíduo é $63\%$ de $\max|\mathbf{u}|$, e o obstáculo deixou de
ser obstáculo.

\paragraph{Advertência de método.} A primeira versão deste varrimento foi até
$2\times10^{-2}$ e \textbf{não encontrou limite algum}. Uma tabela inteiramente
verde \emph{parece} resposta. Varrimento que não encontra a fronteira não mediu o
que se pensa que mediu.

\subsection{O erro do forçamento defasado escala com o número de Courant}

\pgfplotstabletypeset[
  columns/0/.style={column name={$h$}, fixed, precision=4},
  columns/1/.style={column name={cél./$D$}, fixed, precision=0},
  columns/2/.style={column name={$\dt$}, sci, precision=2},
  columns/3/.style={column name={resíduo}, sci, precision=3},
  every head row/.style={before row=\toprule, after row=\midrule},
  every last row/.style={after row=\bottomrule},
]{dados/refino-defasado.dat}

\medskip
Tempo físico fixo ($t = 0{,}2$). O resíduo é $O(\dt)$ \emph{puro} nas duas
malhas --- razões $2{,}04$ e $2{,}02$ --- e \textbf{sem piso}. E a malha fina dá
o dobro da grossa no mesmo $\dt$.

Juntando as duas observações, o resíduo escala com $\dt/h$: o número de Courant.
Faz sentido --- o erro da defasagem é $\dt\,\partial_t u$, e junto ao corpo
$\partial_t u \sim u^2/h$. \textbf{Refinar a malha sem corrigir a defasagem
piora o resíduo}, e é por isso que refinamento local só faz sentido \emph{depois}
da mudança de formulação.

\subsection{Rota pós-preditor e forçamento iterado}

\pgfplotstabletypeset[
  columns/0/.style={column name={iterações}, fixed, precision=0},
  columns/1/.style={column name={resíduo}, sci, precision=3},
  columns/2/.style={column name={ganho vs.\ defasado}, fixed, precision=1},
  every head row/.style={before row=\toprule, after row=\midrule},
  every last row/.style={after row=\bottomrule},
]{dados/iteracoes.dat}

\medskip
Com $\dt = 10^{-3}$ e $t = 0{,}2$. Uma aplicação corta o resíduo pela
\emph{metade}, não o zera --- consequência direta de
$\mathcal{S}\mathcal{I}\neq\mathcal{I}$, e a razão pela qual a literatura
itera~\cite{wang2008}. A convergência satura: dobrar de 5 para 10 iterações
compra $1{,}38\times$; de 10 para 20, $1{,}40\times$.

\subsection{\texorpdfstring{$\dt$}{dt} com a rota pós-preditor}

\pgfplotstabletypeset[
  columns/0/.style={column name={$\dt$}, sci, precision=2},
  columns/1/.style={column name={Uhlmann, 5 it.}, sci, precision=3},
  columns/2/.style={column name={defasado}, sci, precision=3},
  every head row/.style={before row=\toprule, after row=\midrule},
  every last row/.style={after row=\bottomrule},
]{dados/varre-dt-uhlmann.dat}

\medskip
Abaixo de $2\times10^{-3}$ o resíduo \textbf{acha um patamar} ($\approx
1{,}7\times10^{-3}$): é o piso que a formulação defasada nunca deixou aparecer,
porque a rampa $O(\dt)$ o encobria.

\paragraph{O ganho se converte em custo de máquina.} A rota pós-preditor em
$\dt = 2\times10^{-2}$ é tão acurada quanto a defasada em $\dt = 10^{-3}$:
\textbf{vinte vezes o passo para a mesma acurácia}.

\subsection{Refinamento: \texorpdfstring{$\Cd$}{Cd} contra a referência}

\pgfplotstabletypeset[
  columns/0/.style={column name={cél./$D$}, fixed, precision=0},
  columns/1/.style={column name={$h$}, fixed, precision=4},
  columns/2/.style={column name={$\Cd$}, fixed, precision=6},
  columns/3/.style={column name={$\Cl$}, fixed, precision=6},
  columns/4/.style={column name={resíduo}, sci, precision=3},
  every head row/.style={before row=\toprule, after row=\midrule},
  every last row/.style={after row=\bottomrule},
]{dados/refino-cd.dat}

\medskip
Rota pós-preditor com 5 iterações, $\dt = 5\times10^{-3}$, medidos no
\emph{mesmo} instante $t = 14$, com $\Cd$ estável na quinta casa decimal entre
quadros consecutivos.

\begin{center}
\begin{tabular}{lSS}
  \toprule
  & {20 cél./$D$} & {40 cél./$D$} \\
  \midrule
  erro em $\Cd$ & {$+2{,}62\%$} & {$+0{,}41\%$} \\
  \bottomrule
\end{tabular}
\end{center}

A razão entre os erros é $6{,}4$ para um refinamento de $2$, o que dá inclinação
de convergência $\approx 2{,}7$.

\paragraph{Esse valor merece desconfiança.} A literatura mede inclinação
$\approx 1{,}0$ para forçamento direto simples e $\approx 1{,}5$ para
\emph{multi-direct forcing}. Obter $2{,}7$ é quase o dobro da melhor taxa
publicada. Três explicações competem: (i) as iterações compram ordem; (ii) o erro
da malha grossa tinha componentes que se cancelavam; (iii) algum detalhe da
medida de $\Cd$ escala com $h$ de modo não previsto. \textbf{Dois pontos não
distinguem as três} --- daí a terceira resolução.

\paragraph{Observação que separa duas coisas.} A malha fina tem resíduo de não
escorregamento \emph{pior} ($8{,}3\times10^{-3}$ contra $4{,}0\times10^{-3}$) e
$\Cd$ \emph{melhor}. Resíduo e erro no arrasto não são a mesma quantidade: o
primeiro mede escorregamento nos marcadores; o segundo depende de quão bem o
escoamento em torno do corpo está resolvido.

\section{Discussão e trabalho futuro}
\label{sec:discussao}

\subsection{O que este trabalho estabelece}

\begin{itemize}
  \item Um método de fronteira imersa com forçamento direto, verificado nos seus
        operadores de transferência independentemente do solver, em $np=1,2,3$.
  \item Que a formulação defasada tem erro $O(\dt/h)$ --- número de Courant ---
        e que refinar a malha sem corrigi-la \emph{piora} o não escorregamento.
  \item Que a rota pós-preditor com forçamento iterado reduz o resíduo em
        $41\times$ e permite $\dt$ vinte vezes maior para a mesma acurácia.
  \item $\Cd$ com erro de $+2{,}6\%$ a 20 células por diâmetro e $+0{,}4\%$ a 40,
        contra o valor de referência do \emph{benchmark}.
\end{itemize}

\subsection{O que este trabalho não estabelece}

\begin{itemize}
  \item \textbf{A inclinação de convergência.} Dois pontos não a determinam, e o
        valor obtido é quase o dobro do publicado. A terceira resolução está em
        curso.
  \item \textbf{Nada sobre $\Cl$.} O suporte do núcleo é maior que a assimetria
        que o produz.
  \item \textbf{Nada em 3D quantitativamente.} A geometria 3D-1Z existe e roda,
        mas a 10 células por diâmetro e sem comparação com referência.
  \item \textbf{Nada sobre corpo móvel ou deformável.} Ver abaixo.
  \item \textbf{O termo de inércia do fluido interno.} O $\Cd$ calculado ignora a
        inércia do fluido ``fantasma'' dentro do corpo, que Uhlmann contabiliza à
        parte. Para corpo fixo em regime permanente o termo se anula --- e o
        escoamento estacionou ---, de modo que o número é legítimo aqui; num
        transiente, não seria.
\end{itemize}

\subsection{Contorno rígido e interface entre fluidos são problemas distintos}

``Fronteira imersa'' cobre duas linhagens que compartilham a maquinaria e
divergem em quase tudo o mais:

\begin{center}
\begin{tabular}{lll}
  \toprule
   & \textbf{contorno rígido} & \textbf{interface entre fluidos} \\
  \midrule
  força        & multiplicador de Lagrange & constitutiva ($\sigma\kappa\mathbf{n}$) \\
  marcadores   & fixos                     & advectados \\
  conectividade& dispensável               & \textbf{indispensável} (curvatura) \\
  remalhamento & não                       & sim, ou a interface fica permeável \\
  \bottomrule
\end{tabular}
\end{center}

Os operadores de transferência servem aos dois sem diferença. A distribuição por
posse euleriana, porém, \textbf{destrói a ordem da curva} --- correto e barato
para o caso rígido, inviável para o outro, que precisa de vizinhos para calcular
curvatura.

\subsection{Refinamento local}

É o caminho da literatura: o erro é de ordem~1 \emph{junto} à interface e ordem~2
longe dela~\cite{lai2000,griffith2005}, e o método adaptativo de
Roma--Peskin--Berger~\cite{roma1999} --- fonte do núcleo aqui usado --- existe
exatamente para isso. Em 3D é a única rota viável: 40 células por diâmetro
uniformes custariam $\sim$27 milhões de células contra $\sim$420 mil hoje.

Três ressalvas, na ordem em que doem:

\begin{enumerate}
  \item A enumeração do suporte pressupõe um $h$ único e global. Numa malha
        graduada isso não é ajuste, é reescrita.
  \item O núcleo na fronteira de refinamento não preserva de graça a partição da
        unidade nem a propriedade adjunta --- que são precisamente as duas
        propriedades que os testes verificam.
  \item \textbf{Refinamento local reduz células, não passos.} O $\dt$ é global e
        ditado pela menor célula. O custo de refinar é $(1/h)^{D+1}$, e o
        refinamento local ataca apenas $D$ desse expoente.
\end{enumerate}

\subsection{Próximos passos}

\begin{enumerate}
  \item Terceira resolução (80 células/$D$), para decidir a inclinação de
        convergência. \emph{Em curso.}
  \item Se a inclinação se confirmar, refinamento local, com os pré-requisitos
        acima.
  \item Verificação 3D quantitativa contra o \emph{benchmark} 3D-1Z.
  \item Comparação com a rota de Lai--Peskin (ponto médio, duas avaliações de
        força por passo, sem iterar), que é a alternativa estrutural ao
        forçamento iterado.
\end{enumerate}


\bibliographystyle{plain}
\bibliography{referencias}

\end{document}
